\documentclass{article}
\usepackage[utf8]{inputenc}

\usepackage{amsmath, amssymb, amsthm, amsfonts}
\usepackage{subcaption, enumerate}
\usepackage{bbold}
\usepackage{polynom}
\usepackage{xfrac}
\usepackage{graphicx}
\usepackage{mathtools}
\usepackage{enumitem}

\usepackage{tikz}

\theoremstyle{definition}
\newtheorem{definition}{Definition}
\newtheorem{example}{Example}[definition]
\newcommand{\dab}
{
	\begin{tikzpicture}[scale=0.04]
		% head
		\draw[thick] (0, 0) circle (3cm);

		% bent arm
		\draw[thick, rotate around={15 : (0, 0)}] (-3, 0) -- (-3.7, -2) -- (4.1, -1.1);

		% extended arm
		\draw[thick, rotate=20] (3, 0) -- (7, 2);
	\end{tikzpicture}
}
\newcommand{\st}{\text{ s.t. }}

\makeatletter
\renewcommand*\env@matrix[1][*\c@MaxMatrixCols c]{%
  \hskip -\arraycolsep
  \let\@ifnextchar\new@ifnextchar
  \array{#1}}
\makeatother

\usepackage{fullpage}

\title{Linear Transforms Day 5 \\ with Dr Nevard}
\author{Notes by Laura Gallo}
\date{October 1, 2020}

\begin{document}
\maketitle

\section{Gaussian Elimination}
Consider the following system of equations:
\begin{align*}
	x+2y-4z+2w&=5 \\
	2x+y+3z-w&=6
\end{align*}
\\
We can create an \emph{augmented matrix} for the above system and triangulise it:
\begin{equation*}
	\left(
		\begin{matrix}[cccc|c]
			1 & 2 & -4 & 2 & 5 \\
			2 & 1 & 3 & -1 & 6
		\end{matrix}
	\right)
\end{equation*}
\begin{equation*}
	\left(
		\begin{matrix}[cccc|c]
			1 & 2 & -4 & 2 & 5 \\
			0 & -3 & 11 & -5 & -4
		\end{matrix}
	\right)
\end{equation*}
\begin{equation*}
	\left(
		\begin{matrix}[cccc|c]
			1 & 2 & -4 & 2 & 5 \\
			0 & 1 & -\frac{11}{3} & \frac{5}{3} & \frac{4}{3}
		\end{matrix}
	\right)
\end{equation*}
\begin{equation*}
	\left(
		\begin{matrix}[cccc|c]
			1 & 0 & \frac{10}{3} & \frac{4}{3} & \frac{7}{3} \\[1mm]
			0 & 1 & -\frac{11}{3} & \frac{5}{3} & \frac{4}{3}
		\end{matrix}
	\right)
\end{equation*}
\\
From this we can get the following:
\begin{align*}
	x&=\frac{7}{3}-\frac{10}{3}z+\frac{4}{3}w \\
	y&=-\frac{4}{3}+\frac{11}{3}z-\frac{5}{3}w
\end{align*}
\\
This process of solving for $x$ and $y$ using row operations is called Gaussian elimination. Alternatively, we can write the solution set as follows:
\begin{equation*}
	\left(
		\begin{matrix}
			x \\
			y
		\end{matrix}
	\right) = \left(
		\begin{matrix}
			\frac{7}{3} \\
			-\frac{4}{3}
		\end{matrix}
	\right) + \left(
		\begin{matrix}
			-10 \\
			11
		\end{matrix}
	\right)z + \left(
		\begin{matrix}
			4 \\
			-5
		\end{matrix}
	\right)w
\end{equation*}
\\
Recall that \begin{math}
	\left\{\left(
		\begin{matrix}
			-10 \\
			11
		\end{matrix}
	\right),
	\left(
		\begin{matrix}
			4 \\
			-5
		\end{matrix}
	\right)\right\}
\end{math} must span the null space of the linear transform represented by this matrix.

\section{Row Operations as Matrix Multiplication}
Every row operation can be represented as a left matrix multiplication.

\subsection{Swapping rows}
Swapping rows is represented by a special matrix called a \emph{permutation matrix}.

\begin{definition}
	A \emph{permutation matrix} is a matrix that has a 1 in each row and column and a total of $n$ nonzero elements.
\end{definition}

\begin{example}
	\begin{math}
		\left(
			\begin{matrix}
				0 & 1 \\
				1 & 0
			\end{matrix}
		\right)
	\end{math} is a permutation matrix, and is applied as:
	\begin{equation*}
		\left(
			\begin{matrix}
				0 & 1 \\
				1 & 0
			\end{matrix}
		\right) \left(
			\begin{matrix}
				a & b \\
				c & d
			\end{matrix}
		\right) = \left(
			\begin{matrix}
				c & d \\
				a & b
			\end{matrix}
		\right)
	\end{equation*}
\end{example}

\subsection{Subtracting a Multiple of a Row}
Subtracting a multiple of a row is represented by an identity matrix with one of its columns filled with nonzero numbers below the 1. This is an example of a \emph{lower triangular matrix}. An example is shown below:

\begin{equation*}
	\left(
		\begin{matrix}
			1 & 0 & 0 \\
			-x & 1 & 0 \\
			-y & 0 & 1
		\end{matrix}
	\right) \left(
		\begin{matrix}
			a & b & c \\
			d & e & f \\
			g & h & i
		\end{matrix}
	\right) = \left(
		\begin{matrix}
			a & b & c \\
			d-ax & e-bx & f-cx \\
			g-ay & h-by & i-cy
		\end{matrix}
	\right)
\end{equation*}

\subsection{Generalisation}
It can be shown that the following holds for any matrix $A$ with some permutation matrix $P$:
\begin{equation*}
	\left(
		\begin{matrix}
			1 & 0 & 0 & \dots & 0 & 0 \\
			\ell_{11} & 1 & 0 & \dots & 0 & 0 \\
			\ell_{21} & 0 & 1 & \dots & 0 & 0 \\
			\vdots & \vdots & \vdots & \ddots & \vdots & \vdots \\
			\ell_{(m-1)1} & 0 & 0 & \dots & 1 & 0 \\
			\ell_{m1} & 0 & 0 & \dots & 0 & 1 \\
		\end{matrix}
	\right)
	\left(
		\begin{matrix}
			1 & 0 & 0 & \dots & 0 & 0 \\
			0 & 1 & 0 & \dots & 0 & 0 \\
			0 & \ell_{12} & 1 & \dots & 0 & 0 \\
			\vdots & \vdots & \vdots & \ddots & \vdots & \vdots \\
			0 & \ell_{(m-2)2} & 0 & \dots & 1 & 0 \\
			0 & \ell_{(m-1)2} & 0 & \dots & 0 & 1 \\
		\end{matrix}
	\right)
\end{equation*}
\begin{equation*}
	\left(
		\begin{matrix}
			1 & 0 & 0 & \dots & 0 & 0 \\
			0 & 1 & 0 & \dots & 0 & 0 \\
			0 & 0 & 1 & \dots & 0 & 0 \\
			\vdots & \vdots & \vdots & \ddots & \vdots & \vdots \\
			0 & 0 & \ell_{(m-3)3} & \dots & 1 & 0 \\
			0 & 0 & \ell_{(m-2)3} & \dots & 0 & 1 \\
		\end{matrix}
	\right) \cdots
	\left(
		\begin{matrix}
			1 & 0 & 0 & \dots & 0 & 0 \\
			0 & 1 & 0 & \dots & 0 & 0 \\
			0 & 0 & 1 & \dots & 0 & 0 \\
			\vdots & \vdots & \vdots & \ddots & \vdots & \vdots \\
			0 & 0 & \ell_{(m-3)3} & \dots & 1 & 0 \\
			0 & 0 & \ell_{(m-2)3} & \dots & 0 & 1 \\
		\end{matrix}
	\right)PA=\text{an upper triangular matrix}
\end{equation*}

\end{document}
